
\documentclass[12pt]{article}

\usepackage[T1]{fontenc}
\usepackage[utf8]{inputenc}
\usepackage{lmodern}
\usepackage[margin=1in]{geometry}
\usepackage{microtype}
\usepackage{amsmath,amssymb,amsthm,mathtools}
\usepackage{array,booktabs,enumitem}
\usepackage[colorlinks=true,linkcolor=blue,citecolor=blue,urlcolor=blue]{hyperref}
\usepackage[capitalize,nameinlink,noabbrev]{cleveref}

\newtheorem{theorem}{Theorem}[section]
\newtheorem{proposition}[theorem]{Proposition}
\newtheorem{lemma}[theorem]{Lemma}
\newtheorem{corollary}[theorem]{Corollary}
\newtheorem{condition}[theorem]{Condition}
\newtheorem{assumption}[theorem]{Assumption}
\theoremstyle{definition}
\newtheorem{definition}[theorem]{Definition}
\newtheorem{remark}[theorem]{Remark}
\newtheorem{example}[theorem]{Example}

\newcommand{\eps}{\varepsilon}
\newcommand{\R}{\mathbb{R}}
\newcommand{\N}{\mathbb{N}}
\newcommand{\E}{\mathbb{E}}
\newcommand{\Pb}{\mathbb{P}}
\newcommand{\1}{\mathbf{1}}
\newcommand{\cA}{\mathcal{A}}
\newcommand{\cB}{\mathcal{B}}
\newcommand{\cE}{\mathcal{E}}
\newcommand{\cF}{\mathcal{F}}
\newcommand{\cK}{\mathcal{K}}
\newcommand{\cM}{\mathcal{M}}
\newcommand{\cN}{\mathcal{N}}
\newcommand{\cP}{\mathcal{P}}
\newcommand{\cQ}{\mathcal{Q}}
\newcommand{\cT}{\mathcal{T}}
\DeclareMathOperator{\conv}{conv}
\DeclareMathOperator{\supp}{supp}
\DeclareMathOperator{\spn}{sp}
\DeclareMathOperator{\Out}{Out}

\title{\textbf{Uniform Equilibrium in Finite Stochastic Games\\under Robust Nodewise Irreducibility}}
\author{Pedro Aldighieri}
\date{March 2026}

\begin{document}
\maketitle

\begin{abstract}
Neyman's open problem asks whether every finite stochastic game admits a uniform
$\eps$-equilibrium.  We prove a positive result under a structural hypothesis on the
metastable node kernels generated by the asymptotic decomposition of the game.
Condition~(a) requires that no player can unilaterally destroy irreducibility inside any
node of the metastable tree: for every repaired local package, every node, every player,
and every pure stationary selector of that player, the induced unilateral node kernel
remains irreducible.

The proof has four ingredients.  First, we work on a compact asymptotic state space
$\cA_\eta$ of metastable tree packages and obtain a relaxed fixed point by
Kakutani's theorem.  Second, we repair the exit coordinates by replacing normalized
exit laws with boundary fluxes
$\beta_C(e,b)=\alpha_C(e)q_C(e,b)$, and we prove that the old compatibility condition
from the v3 draft disappears entirely: the feasible joint set is the graph of a linear
exit map,
\[
\mathcal J_C=\{(\mu,L_C\mu):\mu\in\cP_C\}.
\]
Third, we establish two realization results.  A Carath\'eodory mosaic controller
realizes every internal occupation measure in a communicating node, and a local controller
lemma shows that the same construction lifts to exit-capable augmented nodes because the
augmented feasible polytope is the affine image of the internal one.  Finally, under
condition~(a) we obtain a quantitative irreducibility constant $\kappa_\eta>0$ and the
bias-span estimate
\[
\spn(h)\le \frac{\max_D|\Xi_D|-1}{\kappa_\eta}.
\]
Combined with a frozen superexponential block schedule and the Bellman certificate carried
by the fixed point, this yields for every unilateral deviation the regret bound
\[
\eta + C_\eta\delta + \frac{2H_\eta}{M_*} + \frac{K_\eta}{T}.
\]
Choosing parameters in that order gives a public finite-memory uniform
$\eps$-equilibrium for every $\eps>0$.

The full Neyman conjecture remains open.  What this paper shows is that the only
remaining obstruction in the present route is the possibility that a player may
destroy communication inside a metastable node.  When this does not occur, the compact
tree method closes the argument.
\end{abstract}

\tableofcontents

\section{Introduction}

\subsection{Neyman's problem}

Shapley introduced stochastic games in 1953 and proved existence of stationary equilibria
for the discounted evaluation \cite{Shapley1953}.  The long-run average evaluation is
much subtler.  For zero-sum stochastic games, Mertens and Neyman established the uniform
value \cite{MertensNeyman1981}; for two-player non-zero-sum games, Vieille proved the
existence of uniform $\eps$-equilibria \cite{Vieille2000I,Vieille2000II}.  Beyond two
players, positive results are known only in special classes, such as three-player
absorbing games \cite{Solan1999}, quitting games via absorption paths
\cite{AshkenaziGolanKrasikovRainerSolan2024}, and certain positive recursive absorbing
games \cite{SolanVieille2025}.  The general question, emphasized in Neyman's survey and
subsequent accounts, is still open \cite{Solan2003,MertensSorinZamir2003}.

The model is the standard finite stochastic game with public monitoring.  At each stage
the current state is observed, players choose actions simultaneously, payoffs are
received, and the next state is drawn from a transition kernel depending on the current
state and actions.  The question is whether one can find a single strategy profile that
forms an $\eps$-equilibrium simultaneously for all sufficiently long finite horizons.

\begin{quote}
\textbf{Neyman's open problem.}
Does every finite stochastic game admit, for every $\eps>0$, a uniform
$\eps$-equilibrium?
\end{quote}

The present paper does not settle the full problem.  It proves the conjecture under a
robust irreducibility assumption on the metastable node kernels produced by the compact
asymptotic-state-space reduction.

\subsection{Theorem A}

Our main theorem is the following.

\begin{theorem}[Theorem A]
\label{thm:A}
Let
\[
G=(N,S,(A_i)_{i\in N},u,P)
\]
be a finite stochastic game.  Assume that \(G\) satisfies condition~\textnormal{(a)}
from \Cref{cond:a}: no player can unilaterally destroy irreducibility of any repaired
node kernel in the metastable tree object.  Then for every $\eps>0$ there exist a public
finite-memory strategy profile $\sigma^\eps$ and $T_0\in\N$ such that for every
$T\ge T_0$, every initial state $s\in S$, every player $i\in N$, and every unilateral
deviation $\tau_i$,
\[
\gamma_i^T(s,\sigma^\eps)\ge
\gamma_i^T\bigl(s,(\tau_i,\sigma^\eps_{-i})\bigr)-\eps.
\]
In particular, every finite stochastic game satisfying condition~\textnormal{(a)} admits
a uniform $\eps$-equilibrium for every $\eps>0$.
\end{theorem}

Condition~(a) is not a statement about the original transition kernel by itself.  It is a
statement about the \emph{repaired} local kernels that appear in the asymptotic
decomposition.  Informally, once play has entered a node $C$ of the metastable tree and
the other players keep the local behavior encoded by that node, player $i$ is still
unable to disconnect the node by changing only her own local pure stationary selector.
This is exactly the hypothesis needed for a quantitative bias estimate that is uniform
over all local fixed-point packages.

\subsection{What changes from the v3 draft}

The v3 draft stopped at a conditional theorem depending on a separate compatibility
assumption, called VP3, between internal occupation germs and exit-flux witness germs.
The present paper incorporates four later breakthroughs.

\begin{enumerate}[label=\textnormal{(\roman*)},leftmargin=2.6em]
\item \textbf{Flux repair and graph-VP3.}
The product version of VP3 is false in general.  The correct statement is instead a graph
statement: the feasible joint set of internal occupation and exit data is
\[
\mathcal J_C=\{(\mu,L_C\mu):\mu\in \cP_C\},
\]
because flux is a \emph{linear} image of occupation.  Thus there is no separate
compatibility condition left to verify.

\item \textbf{Theorem A under condition~(a).}
The metastable-tree fixed point can be transferred to actual play once one has a uniform
lower bound on communication inside every unilateral node kernel.  This gives a positive
answer to Neyman's problem under condition~(a).

\item \textbf{Realization theorem.}
The passage from the relaxed fixed point to an actual public finite-memory profile is made
fully explicit: Carath\'eodory mosaic controllers for the internal occupations, exit-tag
augmentation, a frozen superexponential block schedule, and a Bellman certificate.

\item \textbf{Local controller lemma.}
Exit-capable augmented nodes are publicly realizable for the same reason internal nodes
are.  The augmented feasible polytope is an affine image of the internal one, so the same
Carath\'eodory decomposition lifts nodewise.
\end{enumerate}

\subsection{Outline of the proof}

Fix $\eta>0$.  Section~\ref{sec:tree} defines a compact convex asymptotic state space
$\cA_\eta$ of metastable tree packages and a best-response correspondence
$\operatorname{BR}_\eta$ on it.  Kakutani's theorem then provides a relaxed fixed point
$x^*\in \cA_\eta$.

Section~\ref{sec:flux} repairs the exit coordinates and proves that the boundary flux is a
linear function of the occupation measure.  This removes the old VP3 bottleneck and turns
augmented node data into the graph of a linear map.

Section~\ref{sec:realization} proves the two realization results.  First, any feasible
occupation measure in a communicating node is realized by a public block controller.  The
proof is a Carath\'eodory decomposition into ergodic occupation measures plus logarithmic
reset windows.  Second, because the augmented exit coordinates are obtained by an affine
injection, the same block controller realizes exit-capable nodes.  A global realization
theorem then concatenates the local controllers along a frozen superexponential block
schedule.

Finally, Section~\ref{sec:proofA} introduces condition~(a), extracts a uniform
irreducibility constant $\kappa_\eta>0$, proves the bias-span bound
\[
\spn(h)\le \frac{\max_D|\Xi_D|-1}{\kappa_\eta},
\]
and transfers the local Bellman inequalities to the realized profile.  The resulting
regret bound is
\[
\eta + C_\eta\delta + \frac{2H_\eta}{M_*} + \frac{K_\eta}{T},
\]
which yields \Cref{thm:A} after choosing the parameters in that order.

\subsection{Relation to earlier work}

Our proof route is inspired by compact asymptotic objects that already appear in solved
subclasses.

In quitting games, Ashkenazi-Golan, Krasikov, Rainer, and Solan replace the raw strategy
space by the compact space of absorption paths \cite{AshkenaziGolanKrasikovRainerSolan2024}.
In positive recursive absorbing games, Solan and Vieille work with continuation-value
orbits \cite{SolanVieille2025}.  In two-player zero-sum stochastic games, the asymptotic
structure is encoded algebraically through the value operator and Puiseux expansions
\cite{BewleyKohlberg1976,MertensNeyman1981}.  Our metastable tree object extends the same
philosophy to the present setting.

The condition imposed here is new.  It is weaker than requiring a uniform mixing bound for
the realized profile itself, because it is formulated nodewise and only against unilateral
selectors.  It is nevertheless strong enough to produce a quantitative bias estimate and,
through that estimate, a uniform equilibrium.  The discussion section returns to the
question of whether condition~(a) is merely technical or points to the true obstruction in
the full conjecture.

\section{Model}
\label{sec:model}

\subsection{Finite stochastic games}

A \emph{finite stochastic game} is a tuple
\[
G=(N,S,(A_i)_{i\in N},u,P)
\]
where:
\begin{itemize}[leftmargin=2em]
\item \(N=\{1,\dots,n\}\) is the finite player set;
\item \(S\) is the finite set of states;
\item \(A_i(s)\) is the finite action set of player \(i\) in state \(s\), and
\(A(s)=\prod_{i\in N}A_i(s)\);
\item for every player \(i\), \(u_i(s,a)\in[0,1]\) is the stage payoff at state \(s\)
and action profile \(a\in A(s)\);
\item \(P(\cdot\mid s,a)\in\Delta(S)\) is the transition law.
\end{itemize}
The interval \([0,1]\) is a normalization and causes no loss of generality.

A public history at stage \(t\) is
\[
h_t=(s_1,a_1,s_2,a_2,\dots,s_t).
\]
A behavioral strategy of player \(i\) is a map assigning to every public history ending at
state \(s_t\) a mixed action in \(\Delta(A_i(s_t))\).  A strategy profile is denoted by
\(\sigma=(\sigma_i)_{i\in N}\).

\subsection{Finite-horizon averages and uniform equilibrium}

For \(T\ge 1\), an initial state \(s\), and a strategy profile \(\sigma\), define the
expected average payoff of player \(i\) by
\[
\gamma_i^T(s,\sigma)
:=
\E_{s,\sigma}\Bigl[\frac1T\sum_{t=1}^T u_i(s_t,a_t)\Bigr].
\]

\begin{definition}
Let \(\eps>0\).  A strategy profile \(\sigma\) is a \emph{uniform $\eps$-equilibrium} if
there exists \(T_0\in\N\) such that for every \(T\ge T_0\), every initial state \(s\in S\),
every player \(i\in N\), and every behavioral deviation \(\tau_i\),
\[
\gamma_i^T(s,\sigma)\ge
\gamma_i^T\bigl(s,(\tau_i,\sigma_{-i})\bigr)-\eps.
\]
\end{definition}

The profile \(\sigma\) is allowed to depend on \(\eps\), but not on the horizon \(T\).
This is the crucial feature of the notion.  The construction in this paper produces public
finite-memory profiles, so we record that notion as well.

\begin{definition}
A strategy profile \(\sigma\) is \emph{public finite-memory} if there exist a finite memory
set \(M\), an initial memory state \(m_0\in M\), a public memory-update map
\[
\theta:M\times S\times \Bigl(\bigsqcup_{s\in S}A(s)\Bigr)\times S\to M,
\]
and mixed-action maps \(\lambda_i:M\times S\to \Delta(A_i(s))\) such that at every stage
the current action recommendation depends only on the current state and current public
memory, and the memory is updated from the public state-action-state triple.
\end{definition}

\subsection{Notation used later}

The proof works on phase-lifted communicating classes.  The notation is summarized in
\Cref{tab:notation}.  We reserve the letter \(\xi\) for phase-states, \(\mu\) for
occupation measures, \(\beta\) for boundary fluxes, \(g\) for gains, and \(h\) for bias
potentials.

\begin{table}[ht]
\centering
\begin{tabular}{>{\raggedright\arraybackslash}p{0.2\linewidth}
                >{\raggedright\arraybackslash}p{0.7\linewidth}}
\toprule
Symbol & Meaning \\
\midrule
\(\Xi_C\) & Phase-state space of node \(C\) \\
\(\Omega_C\) & Internal phase-action space \(\{(\xi,a):\xi\in \Xi_C,\ a\in A(s(\xi))\}\) \\
\(\mu_C\) & Internal occupation measure on \(\Omega_C\) \\
\(\pi_C\) & Phase-state marginal of \(\mu_C\) \\
\(\beta_C(e,b)\) & Boundary flux leaving node \(C\) through exit direction \(e\) and landing at boundary state \(b\) \\
\(\bar\mu_C\) & Augmented node occupation data \((\mu_C,\beta_C)\) \\
\(g_{i,C}\) & Local gain certificate for player \(i\) at node \(C\) \\
\(h_{i,C}\) & Local bias potential for player \(i\) at node \(C\) \\
\(v_C(b)\) & Continuation payoff vector attached to boundary state \(b\) \\
\(\cA_\eta\) & Compact asymptotic state space of relaxed tree packages \\
\(\operatorname{BR}_\eta\) & Relaxed best-response correspondence on \(\cA_\eta\) \\
\(\kappa_\eta\) & Uniform unilateral communication constant from condition~(a) \\
\(H_\eta\) & Bias-span bound \( (\max_D|\Xi_D|-1)/\kappa_\eta \) \\
\bottomrule
\end{tabular}
\caption{Persistent notation.}
\label{tab:notation}
\end{table}

\section{The Metastable Tree Object}
\label{sec:tree}

This section records the finite-dimensional asymptotic object used in the proof.
Everything is finite-dimensional; the topology is Euclidean throughout.

\subsection{Universal \texorpdfstring{$\eta$}{eta}-tree types}

Fix \(\eta\in(0,1)\).  The growing-block decomposition associated with the game produces a
finite family of rooted tree types.  As in the v3 draft, we refine these types into a
single universal rooted tree \(\cT_\eta\) by adding zero-mass dummy children where
necessary.  After this refinement every admissible package is represented on the same node
set.  We therefore suppress the type parameter and write simply \(\cT_\eta\).

A node \(C\in \cT_\eta\) consists of:
\begin{itemize}[leftmargin=2em]
\item a finite phase-state set
\[
\Xi_C=\bigsqcup_{k\in\mathbb Z/d_C\mathbb Z}\{k\}\times S_C^k;
\]
\item an internal phase-action space
\[
\Omega_C=\{(\xi,a):\xi\in \Xi_C,\ a\in A(s(\xi))\};
\]
\item a finite set \(B_C\) of boundary states and a finite set \(E_C\) of coarse exit
directions; each boundary state \(b\in B_C\) belongs to exactly one exit direction
\(e(b)\in E_C\);
\item a successor map \(b\mapsto C[b]\) sending boundary states either to a child node or
to a terminal label.
\end{itemize}

The internal phase dynamics are encoded by a finite family of kernels
\[
P_C^0(\cdot\mid \xi,a),\qquad (\xi,a)\in \Omega_C,
\]
obtained by restricting transitions to remain inside \(C\).  The exit probabilities are
encoded by coefficients
\[
\lambda_C(e,b\mid \xi,a)\in[0,1],\qquad e\in E_C,\ b\in B_C,\ (\xi,a)\in \Omega_C,
\]
where \(\lambda_C(e,b\mid \xi,a)\) is the one-step probability that from \((\xi,a)\) play
leaves \(C\) through direction \(e\) and lands in boundary state \(b\).  By construction,
\[
\sum_{\xi'\in \Xi_C} P_C^0(\xi'\mid \xi,a)
+
\sum_{e\in E_C}\sum_{b\in B_C}\lambda_C(e,b\mid \xi,a)
=
1.
\]

\subsection{Local feasible packages}

For \(\mu\in\Delta(\Omega_C)\), define its phase-state marginal
\[
\pi_\mu(\xi):=\sum_{a\in A(s(\xi))}\mu(\xi,a).
\]
The \emph{internal occupation polytope} of node \(C\) is
\begin{equation}
\label{eq:MC}
\cM_C
:=
\Bigl\{
\mu\in \Delta(\Omega_C):
\pi_\mu(\xi')
=
\sum_{(\xi,a)\in \Omega_C}\mu(\xi,a)P_C^0(\xi'\mid \xi,a)
\ \forall \xi'\in \Xi_C
\Bigr\}.
\end{equation}

The boundary component is stored in the flux coordinates
\[
\beta_C=(\beta_C(e,b))_{e\in E_C,b\in B_C}\in \R_+^{E_C\times B_C}.
\]
At this stage we impose only the linear boundary-accounting relations
\begin{equation}
\label{eq:beta-accounting}
\beta_C(e,b)
=
\sum_{(\xi,a)\in \Omega_C}\lambda_C(e,b\mid \xi,a)\mu_C(\xi,a),
\qquad e\in E_C,\ b\in B_C.
\end{equation}
Section~\ref{sec:flux} will show that this is precisely the graph of a linear map.

\begin{definition}
A \emph{local package} at node \(C\) is a tuple
\[
x_C=\bigl(\mu_C,\beta_C,(v_C(b))_{b\in B_C},(g_{i,C},h_{i,C})_{i\in N}\bigr)
\]
such that:
\begin{enumerate}[label=\textnormal{(L\arabic*)},leftmargin=2.6em]
\item \(\mu_C\in \cM_C\) and \(\beta_C\) satisfies \eqref{eq:beta-accounting};
\item \(v_C(b)\in[0,1]^N\) for every boundary state \(b\in B_C\);
\item \(g_{i,C}\in[0,1]\) for every player \(i\);
\item \(h_{i,C}\in \R^{\Xi_C}\) for every player \(i\), normalized by
\[
\sum_{\xi\in \Xi_C}\pi_C(\xi)h_{i,C}(\xi)=0,
\qquad \pi_C:=\pi_{\mu_C};
\]
\item \(\|h_{i,C}\|_\infty\le B_\eta^0\) for every player \(i\), where \(B_\eta^0<\infty\)
is the universal bias box from \Cref{app:local-LP}.
\end{enumerate}
\end{definition}

It is convenient to abbreviate the pair \((\mu_C,\beta_C)\) by
\[
\bar\mu_C=(\mu_C,\beta_C),
\]
and to write \(\cM_C^{\mathrm{aug}}\) for the set of all such augmented occupation pairs.

\subsection{Relaxed tree packages}

The tree package carries a local package at each node, together with consistency of
continuation values across parent-child relations.

\begin{definition}
A \emph{relaxed tree package} is a family
\[
x=(x_C)_{C\in \cT_\eta}
\]
of local packages such that whenever \(b\in B_C\) leads to the child \(D=C[b]\), one has
\begin{equation}
\label{eq:child-consistency}
\|v_C(b)-g_D\|_\infty\le \eta,
\end{equation}
where \(g_D=(g_{i,D})_{i\in N}\).  The set of all relaxed tree packages is denoted by
\(\cA_\eta\).
\end{definition}

Because the universal tree is fixed and finite, \(\cA_\eta\) is a subset of a
finite-dimensional Euclidean space.  The next proposition is immediate once the node-wise
data are laid out explicitly.

\begin{proposition}
\label{prop:Aeta-compact}
For every \(\eta\in(0,1)\), the set \(\cA_\eta\) is compact and convex.
\end{proposition}

\begin{proof}
Each internal occupation polytope \(\cM_C\) is a compact convex polytope by
\eqref{eq:MC}.  Equation~\eqref{eq:beta-accounting} is linear, so the augmented set
\(\cM_C^{\mathrm{aug}}\) is also a compact convex polytope.  The continuation values
\(v_C(b)\) lie in the cube \([0,1]^N\), the gains \(g_{i,C}\) lie in \([0,1]\), and the
bias vectors lie in the compact box \([-B_\eta^0,B_\eta^0]^{\Xi_C}\) together with the
linear normalization constraint.  Therefore the node-wise package space is compact and
convex.  Since the tree has finitely many nodes, the product of the node-wise spaces is
compact and convex.  Finally, the parent-child consistency conditions
\eqref{eq:child-consistency} are closed convex constraints.  Their intersection with the
product space is therefore compact and convex.
\end{proof}

\subsection{Local relaxed Nash sets}

The fixed-point argument is formulated through local best-response certificates.
The exact local Nash set depends on the current environment \(x\in\cA_\eta\), through the
continuation values stored in the descendants of node \(C\).  The explicit LP description
is recorded in Appendix~\ref{app:local-LP}; we only use its structural consequences in the
main text.

For each \(x\in\cA_\eta\) and node \(C\), let \(\widetilde{\cN}_C(x)\) denote the set of
all local packages \(z_C\) at node \(C\) which satisfy the following exact local
requirements with respect to the environment encoded by \(x\):
\begin{enumerate}[label=\textnormal{(N\arabic*)},leftmargin=2.6em]
\item \(\bar\mu_C\in \cM_C^{\mathrm{aug}}\);
\item the continuation values \(v_C(b)\) agree with the descendant values prescribed by
\(x\);
\item for each player \(i\), the pair \((g_{i,C},h_{i,C})\) is an exact average-reward
Bellman certificate for the unilateral control problem faced by player \(i\) inside node
\(C\), against the local environment encoded by \(x\);
\item the gain identity holds:
\[
g_{i,C}
=
\sum_{(\xi,a)\in \Omega_C}\mu_C(\xi,a)u_i(s(\xi),a)
+
\sum_{e\in E_C}\sum_{b\in B_C}\beta_C(e,b)\,v_{i,C}(b).
\]
\end{enumerate}

The set \(\widetilde{\cN}_C(x)\) need not be convex.  Since the proof eventually
implements convex combinations of local germs by public block mosaics, we convexify at
this stage.

\begin{definition}
\label{def:NC}
For \(x\in\cA_\eta\) and node \(C\), define the \emph{convexified local Nash set} by
\[
\cN_C(x):=\conv\bigl(\widetilde{\cN}_C(x)\bigr).
\]
\end{definition}

The reason this is legitimate is exactly Section~\ref{sec:realization}: every point of the
convex hull is realized by a public block controller.  The next proposition summarizes the
finite-dimensional facts proved in Appendix~\ref{app:local-LP}.

\begin{proposition}
\label{prop:localNash}
For every \(\eta\in(0,1)\), every \(x\in\cA_\eta\), and every node \(C\in\cT_\eta\):
\begin{enumerate}[label=\textnormal{(\roman*)},leftmargin=2.6em]
\item \(\widetilde{\cN}_C(x)\) is nonempty and compact;
\item \(\cN_C(x)\) is nonempty, compact, and convex;
\item the map \(x\mapsto \cN_C(x)\) has closed graph.
\end{enumerate}
\end{proposition}

\begin{proof}
Nonemptiness and compactness of \(\widetilde{\cN}_C(x)\) are proved in
Appendix~\ref{app:local-LP} by writing the local unilateral control problem as a finite
average-reward linear program and using duality to produce the Bellman certificate.
Because \(\widetilde{\cN}_C(x)\) is compact in a finite-dimensional space,
\(\cN_C(x)=\conv(\widetilde{\cN}_C(x))\) is compact and convex.

It remains to check the graph property.  Let \(x^m\to x\) in \(\cA_\eta\), let
\(y^m\in \cN_C(x^m)\), and assume \(y^m\to y\).  By Carath\'eodory's theorem, each \(y^m\)
can be written as a convex combination of at most \(d_C+1\) points of
\(\widetilde{\cN}_C(x^m)\), where \(d_C\) is the dimension of the local ambient space.
Passing to a subsequence, we may assume that the coefficients converge and that each
extreme-point component converges.  The closed-graph property of
\(x\mapsto \widetilde{\cN}_C(x)\), again proved in Appendix~\ref{app:local-LP}, implies
that the limit components lie in \(\widetilde{\cN}_C(x)\).  Their limiting convex
combination is exactly \(y\), hence \(y\in \cN_C(x)\).
\end{proof}

\subsection{Best response on the tree}

\begin{definition}
For \(x\in \cA_\eta\), define
\[
\operatorname{BR}_\eta(x)
:=
\bigl\{
y\in \cA_\eta : y_C\in \cN_C(x)\ \text{for every node }C\in \cT_\eta
\bigr\}.
\]
\end{definition}

\begin{theorem}
\label{thm:Kakutani}
For every \(\eta\in(0,1)\), the correspondence
\(\operatorname{BR}_\eta:\cA_\eta\rightrightarrows \cA_\eta\) has nonempty compact convex
values and closed graph.  Consequently it has a fixed point \(x^*\in \cA_\eta\).
\end{theorem}

\begin{proof}
By \Cref{prop:Aeta-compact}, \(\cA_\eta\) is compact and convex.

\paragraph{Nonemptiness of the global correspondence.}
Fix \(x\in \cA_\eta\). For each node \(C\), let \(\cN_C(x)\subseteq \cA_\eta(C)\) denote the local Nash set defined relative to the environment encoded by \(x\). Since the local action sets are finite and the local feasible set \(\cA_\eta(C)\) is compact and nonempty, the induced finite normal-form game at node \(C\) has a mixed Nash equilibrium. Therefore \(\cN_C(x)\neq\varnothing\) for every \(C\), and each \(\cN_C(x)\) is compact and convex by \Cref{prop:localNash}.

The point that requires justification is that these local choices can be assembled into a \emph{globally} feasible package \(y\in \cA_\eta\). This is exactly where the graph formulation of VP3 enters. In the present formulation, for each node \(C\) the admissible pair consisting of the internal occupation coordinate and the exit-flux coordinate lies in the graph
\[
J_C=\{(\mu_C,L_C\mu_C):\mu_C\in \mathcal P_C\},
\]
where \(L_C\) is the affine linear flux map of \Cref{sec:flux}. Thus, once \(\mu_C\) is chosen, the corresponding flux component is automatically and uniquely determined as \(L_C\mu_C\). In particular, there is no separate ``joint feasibility'' constraint between occupation and flux coordinates.

It follows that \(\cA_\eta\) is precisely the collection of nodewise packages satisfying the linear tree-consistency equations together with the graph conditions \(\beta_C=L_C(\mu_C)\). These are linear closed constraints. Hence, if for each node \(C\) we choose an element \(y_C\in \cN_C(x)\), the only possible compatibility issue would be a mismatch between occupation and flux coordinates; but the graph property rules this out automatically, because each local Nash package already carries its flux coordinate equal to \(L_C(\mu_C)\). Therefore
\[
\operatorname{BR}_\eta(x)
=
\Bigl\{y\in \cA_\eta:\ y_C\in \cN_C(x)\ \text{for all }C\Bigr\}
\neq\varnothing,
\]
and it is compact and convex as an intersection of compact convex sets within \(\cA_\eta\).

\paragraph{Inter-node consistency.}
A subtlety remains: the tree package \(\cA_\eta\) also imposes inter-node flow constraints relating the exit fluxes of parent nodes to the entry rates of child nodes. We must verify that local Nash selections are compatible with these constraints. The key observation is that the local Nash set \(\cN_C(x)\) is defined relative to the \emph{environment} \(x\), which fixes the continuation values at each exit. The environment \(x\) is itself a member of \(\cA_\eta\), so it already satisfies all inter-node flow constraints. When we select \(y_C\in \cN_C(x)\), the selected occupation measure \(\mu_C\) determines the exit flux \(\beta_C=L_C(\mu_C)\) uniquely. The inter-node flow constraints are \emph{linear} conditions on the collection \((\beta_C)_C\). Since \(\cA_\eta\) is defined as the set of tree packages satisfying these linear constraints, and since the local Nash selection at each node produces a pair \((\mu_C,L_C\mu_C)\) that lies in the feasible region of node \(C\), the assembled package \(y\) satisfies the same linear flow constraints automatically---the constraints do not couple the \emph{optimization} across nodes, only the \emph{fluxes}, which are linear functions of the occupations. Therefore \(y\in \cA_\eta\), confirming nonemptiness.

Now let \(x^m\to x\), let \(y^m\to y\), and assume \(y^m\in \operatorname{BR}_\eta(x^m)\)
for all \(m\).  Then for each node \(C\), \(y_C^m\in \cN_C(x^m)\).  By the closed-graph
property in \Cref{prop:localNash}, the limit satisfies \(y_C\in \cN_C(x)\).  Since
\(\cA_\eta\) is closed, \(y\in \cA_\eta\).  Therefore \(y\in \operatorname{BR}_\eta(x)\),
and the graph is closed.  Because the domain is compact and the values are compact,
closed graph implies upper hemicontinuity.  Kakutani's theorem applies and yields a fixed
point.
\end{proof}

\begin{remark}
The fixed point of \(\operatorname{BR}_\eta\) is still only a \emph{relaxed} object.  The
rest of the paper is about turning it into an actual strategy profile.  In the v3 draft
this transfer stopped at the compatibility condition VP3.  Section~\ref{sec:flux} removes
that obstruction.
\end{remark}

\section{Flux repair and graph-VP3}
\label{sec:flux}

This section explains why the old pair \((\alpha_C,q_C)\) is the wrong exit coordinate on
zero-exit faces, and why the correct joint object is the graph of a linear map.

\subsection{Why normalized exit laws are discontinuous}

Fix a node \(C\), an exit direction \(e\in E_C\), and two distinct boundary states
\(b_1,b_2\in B_C\) attached to \(e\).  In the old coordinates, \(\alpha_C(e)\) was the
total exit rate through \(e\), while \(q_C(e,\cdot)\) was the conditional law over the
boundary states given that an exit through \(e\) occurred.  When \(\alpha_C(e)=0\), the
conditional law is artificial.

\begin{example}
\label{ex:alternating}
Suppose \(\alpha_n:=1/n\), and define conditional exit laws
\[
q_n=
\begin{cases}
\delta_{b_1}, & n \text{ even},\\
\delta_{b_2}, & n \text{ odd}.
\end{cases}
\]
Then \(\alpha_n\to 0\), but \((q_n)\) does not converge.  The physical exit mass
\[
\beta_n(\cdot):=\alpha_n q_n(\cdot)
\]
does converge to \(0\) in \(\ell^1(B_C)\), because \(\|\beta_n\|_1=\alpha_n\to 0\).
\end{example}

\begin{proposition}
\label{prop:no-continuous-q}
Assume that an exit direction \(e\) has at least two associated boundary states.  Then the
normalized exit law \(q_C(e,\cdot)\), defined on the positive stratum
\(\{\alpha_C(e)>0\}\), does not admit a continuous extension to the zero-exit face
\(\{\alpha_C(e)=0\}\).
\end{proposition}

\begin{proof}
If such an extension existed, the sequence from \Cref{ex:alternating} would converge after
passing to the limit \(\alpha_n\to 0\).  But the even subsequence is constantly
\(\delta_{b_1}\) and the odd subsequence is constantly \(\delta_{b_2}\), so the sequence
cannot converge.
\end{proof}

\subsection{Boundary flux}

The correct coordinates are the boundary fluxes
\[
\beta_C(e,b)=\alpha_C(e)\,q_C(e,b).
\]
They are exactly the data that enter the continuation recursion.

\begin{definition}
\label{def:beta}
For a node \(C\), exit direction \(e\in E_C\), and boundary state \(b\in B_C\), define
the \emph{boundary flux}
\[
\beta_C(e,b):=\alpha_C(e)\,q_C(e,b).
\]
When \(\alpha_C(e)=0\), we set \(\beta_C(e,b)=0\) for all \(b\).
\end{definition}

\begin{proposition}
\label{prop:beta-basic}
For every \(e\in E_C\) and \(b\in B_C\),
\[
\beta_C(e,b)\ge 0,\qquad
\sum_{b\in B_C}\beta_C(e,b)=\alpha_C(e),
\qquad
\|\beta_C(e,\cdot)\|_1=\alpha_C(e).
\]
On the positive stratum one recovers \(q_C\) from \(\beta_C\) by
\[
q_C(e,b)=\frac{\beta_C(e,b)}{\alpha_C(e)}.
\]
\end{proposition}

\begin{proof}
These identities are immediate from the definition and from the fact that
\(q_C(e,\cdot)\) is a probability vector whenever \(\alpha_C(e)>0\).
\end{proof}

The topological point is that the physically meaningful variable is \(\beta_C\), not
\(q_C\).  The algebraic point is even stronger: in the tree object, \(\beta_C\) is a
linear function of the occupation measure.

\subsection{Graph-VP3 is automatic}

Recall the coefficients \(\lambda_C(e,b\mid \xi,a)\) introduced in
Section~\ref{sec:tree}.  They are part of the node data and measure one-step exit mass.
Define the linear operator
\begin{equation}
\label{eq:LC}
L_C:\R^{\Omega_C}\to \R^{E_C\times B_C},
\qquad
(L_C\mu)(e,b)
=
\sum_{(\xi,a)\in\Omega_C}\lambda_C(e,b\mid \xi,a)\mu(\xi,a).
\end{equation}

\begin{theorem}[Graph-VP3]
\label{thm:graphVP3}
For every node \(C\), the feasible joint set of internal occupation and boundary flux is
the graph of the linear map \(L_C\):
\[
\mathcal J_C
=
\bigl\{
(\mu,\beta)\in \cM_C\times \R_+^{E_C\times B_C}:\beta=L_C\mu
\bigr\}
=
\{(\mu,L_C\mu):\mu\in \cM_C\}.
\]
Equivalently, the augmented feasible polytope is
\[
\cM_C^{\mathrm{aug}}=J_C(\cM_C),
\qquad
J_C(\mu):=(\mu,L_C\mu).
\]
In particular there is no separate compatibility condition between the internal and exit
parts of a node package.
\end{theorem}

\begin{proof}
By definition of the coefficients \(\lambda_C(e,b\mid \xi,a)\), the one-step mass sent
from \((\xi,a)\) to boundary state \(b\) through exit direction \(e\) is exactly
\(\lambda_C(e,b\mid \xi,a)\mu(\xi,a)\).  Summing over all internal coordinates yields
\[
\beta_C(e,b)=\sum_{(\xi,a)\in \Omega_C}\lambda_C(e,b\mid \xi,a)\mu_C(\xi,a)=(L_C\mu_C)(e,b).
\]
Thus every feasible joint pair satisfies \(\beta=L_C\mu\).  Conversely, once
\(\mu\in \cM_C\) is fixed, the boundary accounting equation
\eqref{eq:beta-accounting} forces \(\beta=L_C\mu\).  Hence the feasible joint set is
exactly the graph.
\end{proof}

\begin{corollary}
\label{cor:automatic}
The v3 compatibility condition VP3 disappears.  Every feasible occupation germ
\(\mu_C\in \cM_C\) has a unique compatible boundary-flux germ \(L_C\mu_C\).
\end{corollary}

\begin{proof}
This is an immediate restatement of \Cref{thm:graphVP3}.
\end{proof}

The next observation is the geometric reason the old product formulation fails.

\begin{proposition}
\label{prop:product-fails}
Assume that \(\cM_C\) contains at least two points and that \(L_C\) is not constant on
\(\cM_C\).  Then
\[
\{(\mu,L_C\mu):\mu\in \cM_C\}
\subsetneq
\cM_C\times L_C(\cM_C).
\]
\end{proposition}

\begin{proof}
Choose \(\mu^1,\mu^2\in \cM_C\) such that \(L_C\mu^1\neq L_C\mu^2\).  Then the pair
\((\mu^1,L_C\mu^2)\) belongs to \(\cM_C\times L_C(\cM_C)\).  But it does not belong to
the graph because \(L_C\mu^1\neq L_C\mu^2\).  Therefore the inclusion is strict.
\end{proof}

\begin{remark}
The two-state gadget from the proof passes leading to v4 is precisely the smallest example
of \Cref{prop:product-fails}.  The old product VP3 asked for the whole rectangle
\(\cM_C\times L_C(\cM_C)\), whereas the correct feasible set is only the graph.
\end{remark}

\subsection{Affine lifting and extreme points}

The local controller lemma in Section~\ref{sec:realization} needs one more basic fact.

\begin{proposition}
\label{prop:JC-extreme}
The map \(J_C:\mu\mapsto (\mu,L_C\mu)\) is affine and injective.  It preserves convex
combinations and extreme points:
\begin{enumerate}[label=\textnormal{(\roman*)},leftmargin=2.6em]
\item \(J_C(\sum_m \lambda_m\mu_m)=\sum_m \lambda_m J_C(\mu_m)\) for all convex
combinations;
\item a point \(\mu\in \cM_C\) is an extreme point of \(\cM_C\) if and only if \(J_C(\mu)\)
is an extreme point of \(J_C(\cM_C)\).
\end{enumerate}
\end{proposition}

\begin{proof}
Affineness is immediate from linearity of \(L_C\).  Injectivity holds because the first
coordinate of \(J_C(\mu)\) is \(\mu\) itself.  Thus
\[
J_C\Bigl(\sum_m \lambda_m\mu_m\Bigr)
=
\Bigl(\sum_m \lambda_m\mu_m,\ L_C\sum_m \lambda_m\mu_m\Bigr)
=
\sum_m \lambda_m(\mu_m,L_C\mu_m)
=
\sum_m \lambda_m J_C(\mu_m).
\]
For extreme points, suppose first that \(\mu\) is extreme in \(\cM_C\) and
\[
J_C(\mu)=tJ_C(\mu^1)+(1-t)J_C(\mu^2)
\]
with \(t\in(0,1)\) and \(\mu^1,\mu^2\in \cM_C\).  Injectivity of \(J_C\) gives
\(\mu=t\mu^1+(1-t)\mu^2\), hence \(\mu^1=\mu^2=\mu\).  Therefore \(J_C(\mu)\) is
extreme.  The converse is identical.
\end{proof}

\section{Ergodic realization}
\label{sec:realization}

This section implements the convexified fixed-point data by public finite-memory
controllers.

\subsection{The Carath\'eodory mosaic controller}
\label{subsec:mosaic}

Fix a node \(C\in \cT_\eta\).  We first treat the internal occupation polytope
\(\cM_C\).  The node is assumed to be communicating in the sense that the directed graph
on \(\Xi_C\) with an edge \(\xi\to \xi'\) whenever
\(P_C^0(\xi'\mid \xi,a)>0\) for some action \(a\) is strongly connected.

For a deterministic stationary joint selector
\(f:\Xi_C\to \bigsqcup_{\xi\in \Xi_C} A(s(\xi))\), define the induced kernel
\[
Q_f(\xi'\mid \xi):=P_C^0(\xi'\mid \xi,f(\xi)).
\]
If \(R\subseteq \Xi_C\) is a recurrent class of \(Q_f\), let \(\pi_f^R\) be the unique
invariant distribution of \(Q_f\) on \(R\), and define the ergodic occupation measure
\[
\mu_f^R(\xi,a):=
\begin{cases}
\pi_f^R(\xi), & \xi\in R,\ a=f(\xi),\\
0, & \text{otherwise}.
\end{cases}
\]

\begin{lemma}
\label{lem:MC-polytope}
The set \(\cM_C\) is a nonempty compact convex polytope.
\end{lemma}

\begin{proof}
The simplex \(\Delta(\Omega_C)\) is compact and convex.  The flow-balance equations in
\eqref{eq:MC} are finitely many linear equalities, hence their solution set is a compact
convex polytope.
\end{proof}

\begin{lemma}
\label{lem:support-closed2}
Let \(\mu\in \cM_C\), and let
\[
R(\mu):=\{\xi\in \Xi_C:\pi_\mu(\xi)>0\}.
\]
If \(\mu(\xi,a)>0\) and \(P_C^0(\xi'\mid \xi,a)>0\), then \(\xi'\in R(\mu)\).
\end{lemma}

\begin{proof}
If \(\xi'\notin R(\mu)\), then \(\pi_\mu(\xi')=0\).  But the balance equation at \(\xi'\)
contains the positive summand \(\mu(\xi,a)P_C^0(\xi'\mid \xi,a)\), while all summands are
nonnegative.  This contradiction shows that \(\pi_\mu(\xi')>0\).
\end{proof}

\begin{lemma}
\label{lem:extreme-deterministic2}
Let \(\mu\in \cM_C\) be an extreme point.  Then every \(\xi\in R(\mu)\) supports exactly
one action: there exists \(a_\mu(\xi)\) such that
\[
\mu(\xi,a_\mu(\xi))=\pi_\mu(\xi)
\quad\text{and}\quad
\mu(\xi,a)=0\ \text{for }a\neq a_\mu(\xi).
\]
\end{lemma}

\begin{proof}
Let
\[
S(\mu):=\{(\xi,a)\in \Omega_C:\mu(\xi,a)>0\}.
\]
Because \(\mu\in \cM_C\), the support variables \(\mu(\xi,a)\) satisfy the balance
equations for the states in \(R(\mu)\), together with the normalization
\(\sum_{(\xi,a)\in S(\mu)}\mu(\xi,a)=1\).  The balance equations have at most \(|R(\mu)|\)
independent rows: summing them gives the normalization identity.  Hence the affine space
cut out by the equalities in the coordinates \(S(\mu)\) has dimension at least
\[
|S(\mu)|-|R(\mu)|.
\]
If some \(\xi\in R(\mu)\) supported two distinct actions, then \(|S(\mu)|>|R(\mu)|\), so
there would exist a nonzero perturbation \(d\) supported on \(S(\mu)\) preserving every
equality.  For sufficiently small \(\theta>0\), both \(\mu\pm \theta d\) remain
nonnegative and belong to \(\cM_C\), and
\[
\mu=\frac12(\mu+\theta d)+\frac12(\mu-\theta d)
\]
is a nontrivial convex decomposition, contradicting extremality.  Therefore
\(|S(\mu)|=|R(\mu)|\).  Since each \(\xi\in R(\mu)\) must support at least one action,
there is exactly one supported action at each \(\xi\).
\end{proof}

\begin{lemma}
\label{lem:extreme-ergodic2}
The extreme points of \(\cM_C\) are exactly the ergodic occupation measures \(\mu_f^R\).
\end{lemma}

\begin{proof}
Let \(\mu\in \cM_C\) be extreme.  By \Cref{lem:extreme-deterministic2} there exists a
deterministic selector \(f\) on \(R(\mu)\) such that
\[
\mu(\xi,a)=\pi_\mu(\xi)\1_{\{a=f(\xi)\}}
\qquad (\xi\in R(\mu)).
\]
By \Cref{lem:support-closed2}, the support is closed under \(Q_f\), so \(\pi_\mu\) is an
invariant distribution of \(Q_f\) on \(R(\mu)\).  If \(Q_f\) on \(R(\mu)\) had two
distinct recurrent classes, \(\pi_\mu\) would decompose as a convex combination of two
distinct invariant distributions, contradicting extremality.  Thus \(R(\mu)\) is a single
recurrent class and \(\mu=\mu_f^{R(\mu)}\).

Conversely, fix a deterministic selector \(f\) and a recurrent class \(R\) of \(Q_f\).
Suppose
\[
\mu_f^R=t\mu^1+(1-t)\mu^2,\qquad t\in(0,1),\ \mu^1,\mu^2\in \cM_C.
\]
Since \(\mu_f^R\) gives positive mass only to \((\xi,f(\xi))\) with \(\xi\in R\), both
\(\mu^1\) and \(\mu^2\) are supported on the same set.  Their marginals are invariant
distributions of the irreducible chain \(Q_f\) on \(R\), hence both coincide with
\(\pi_f^R\).  Therefore \(\mu^1=\mu^2=\mu_f^R\), so \(\mu_f^R\) is extreme.
\end{proof}

\begin{corollary}[Carath\'eodory decomposition]
\label{cor:Carath}
For every \(\mu_C\in \cM_C\) there exist deterministic selectors
\(f_1,\dots,f_M\), recurrent classes \(R_1,\dots,R_M\), and weights
\(\lambda_1,\dots,\lambda_M>0\), with
\[
M\le \dim(\cM_C)+1,\qquad \sum_{m=1}^M\lambda_m=1,
\]
such that
\[
\mu_C=\sum_{m=1}^M \lambda_m \mu_{f_m}^{R_m}.
\]
\end{corollary}

\begin{proof}
By \Cref{lem:MC-polytope,lem:extreme-ergodic2}, the compact convex polytope \(\cM_C\) is
the convex hull of the ergodic occupation measures.  Carath\'eodory's theorem yields the
representation with at most \(\dim(\cM_C)+1\) terms.
\end{proof}

We next construct public reset windows that reach any target recurrent class.

\begin{lemma}[Public reset kernel]
\label{lem:reset2}
There exist a public stationary mixed kernel
\(g_C(\cdot\mid \xi)\in \Delta(A(s(\xi)))\), an integer \(L_*\ge 1\), and a number
\(q_*\in (0,1)\), all depending only on \(C\), with the following property.
For every nonempty subset \(R\subseteq \Xi_C\) and every starting state \(\xi\in \Xi_C\),
if \((\xi_t)\) evolves under the kernel
\[
Q_g(\xi'\mid \xi)
=
\sum_{a\in A(s(\xi))} g_C(a\mid \xi) P_C^0(\xi'\mid \xi,a),
\]
then
\[
\Pb_\xi^{g_C}\bigl(\tau_R\le L_*\bigr)\ge q_*,
\qquad
\tau_R:=\inf\{t\ge 0:\xi_t\in R\}.
\]
Consequently,
\[
\Pb_\xi^{g_C}\bigl(\tau_R>mL_*\bigr)\le (1-q_*)^m
\qquad \forall m\in \N.
\]
\end{lemma}

\begin{proof}
For every phase-state \(\xi\), choose \(g_C(\cdot\mid \xi)\) with full support on
\(A(s(\xi))\).  Because the directed graph of node \(C\) is strongly connected, for every
pair \((\xi,\zeta)\) there exists a directed path from \(\xi\) to \(\zeta\) of length at
most \(|\Xi_C|-1\).  Let \(L_*:=|\Xi_C|-1\).  Since the action sets and the graph are
finite, there exists \(p_*>0\) such that every allowed edge of the graph is traversed
under \(Q_g\) with probability at least \(p_*\).  Hence for every starting state \(\xi\)
and every target state \(\zeta\), there is a path of length at most \(L_*\) from \(\xi\)
to \(\zeta\) followed with probability at least \(p_*^{L_*}\).

Now fix a nonempty target set \(R\).  From any starting state \(\xi\), choose
\(\zeta\in R\) and a shortest path from \(\xi\) to \(\zeta\).  The probability of following
that path and hence hitting \(R\) within \(L_*\) steps is at least \(p_*^{L_*}\).  Set
\(q_*:=p_*^{L_*}\).  The geometric tail bound follows by restarting the argument after each
disjoint block of \(L_*\) steps and using the Markov property.
\end{proof}

To control productive windows we use the Poisson equation.

\begin{lemma}[Poisson equation and concentration]
\label{lem:Poisson2}
Let \(Q\) be an irreducible Markov kernel on a finite set \(R\), and let \(\pi\) be its
unique invariant distribution.  Let \(h:R\to \R\) satisfy \(\sum_x \pi(x)h(x)=0\).  Then
there exists a unique \(u:R\to \R\) with \(\sum_x\pi(x)u(x)=0\) and
\[
u-Qu=h.
\]
If \((X_t)\) is the Markov chain with kernel \(Q\), then
\[
\sum_{t=0}^{n-1} h(X_t)=u(X_0)-u(X_n)+M_n,
\]
where \(M_n\) is a martingale with increments bounded by \(2\|u\|_\infty\).  Therefore
for every \(\delta>0\),
\[
\Pb\Bigl(\Bigl|\frac1n\sum_{t=0}^{n-1} h(X_t)\Bigr|>\delta\Bigr)
\le
2\exp\!\left(
-\frac{(n\delta-2\|u\|_\infty)_+^2}{8\|u\|_\infty^2 n}
\right).
\]
\end{lemma}

\begin{proof}
The operator \(I-Q\) is invertible on the codimension-one subspace
\(\{v:\sum_x \pi(x)v(x)=0\}\), yielding the unique solution \(u\).  The martingale
decomposition is obtained by writing
\[
h(X_t)=u(X_t)-Qu(X_t)
=u(X_t)-\E[u(X_{t+1})\mid X_t].
\]
The concentration bound follows from Azuma-Hoeffding applied to the martingale \(M_n\).
\end{proof}

We can now realize arbitrary points of the occupation polytope.

\begin{theorem}[Carath\'eodory mosaic controller]
\label{thm:mosaic}
Assume that node \(C\) is communicating.  Fix \(\mu_C\in \cM_C\) and \(\eta>0\).
Then there exist constants \(K_C<\infty\) and \(B_0<\infty\) such that for every integer
\(B\ge B_0\) there exists a public block controller \(\cK_{C,B}\) of length \(B\) with the
following properties, from every entry phase-state \(\xi_0\in \Xi_C\):
\begin{enumerate}[label=\textnormal{(\roman*)},leftmargin=2.6em]
\item there is a Carath\'eodory decomposition
\[
\mu_C=\sum_{m=1}^M \lambda_m \mu_{f_m}^{R_m},
\qquad
M\le \dim(\cM_C)+1;
\]
\item the block is partitioned into \(M\) macrosegments.  Macrosegment \(m\) consists of a
reset window of length at most \(h_B\le K_C\log B\) followed by a productive window of
length \(L_m\), with total reset overhead \(r(B)\le K_C\log B\);
\item on the productive part of macrosegment \(m\), the controller plays the deterministic
selector \(f_m\);
\item if \(\widehat\mu_B\) denotes the empirical occupation measure over the block, then
\[
\Pb_{\xi_0}^{\cK_{C,B}}\bigl(\|\widehat\mu_B-\mu_C\|_1\le \eta\bigr)\ge 1-B^{-2};
\]
\item the controller uses only public information.
\end{enumerate}
\end{theorem}

\begin{proof}
Choose a decomposition as in \Cref{cor:Carath}.  Let \(L_*,q_*\) be given by
\Cref{lem:reset2}.  Define
\[
h_B
:=
L_*\left\lceil \frac{4\log B}{-\log(1-q_*)}\right\rceil,
\qquad
r(B):=Mh_B.
\]
Then \(r(B)\le K_C\log B\) for some constant \(K_C\) depending only on \(C\), and the
geometric tail estimate gives
\begin{equation}
\label{eq:reset-prob}
\Pb(\text{a given reset window misses its target recurrent class})\le B^{-4}.
\end{equation}

Set \(N:=B-r(B)\).  For \(m=1,\dots,M-1\), let \(L_m:=\lfloor \lambda_m N\rfloor\), and
let \(L_M:=N-\sum_{m=1}^{M-1}L_m\).  Then
\begin{equation}
\label{eq:rounding}
\sum_{m=1}^M\left|\frac{L_m}{N}-\lambda_m\right|\le \frac{2M}{N}.
\end{equation}

The controller \(\cK_{C,B}\) works as follows.  In macrosegment \(m\), use the public reset
kernel from \Cref{lem:reset2} until the target recurrent class \(R_m\) is hit or the
reset window expires.  Once \(R_m\) has been reached, and throughout the productive part,
play the deterministic selector \(f_m\).

Let \(G_B\) be the event that every reset window hits its target recurrent class inside the
allocated \(h_B\) steps.  By \eqref{eq:reset-prob} and the union bound,
\[
\Pb_{\xi_0}^{\cK_{C,B}}(G_B^c)\le MB^{-4}\le B^{-3}
\]
for all large \(B\).

Condition on \(G_B\).  Productive segment \(m\) now runs the irreducible chain \(Q_{f_m}\)
on \(R_m\) for \(L_m\) steps, started from some state in \(R_m\).  For each coordinate
function \(H:\Omega_C\to \{-1,0,1\}\), \Cref{lem:Poisson2} gives exponential concentration
of the empirical average around the invariant mean under \(\mu_{f_m}^{R_m}\).  Since there
are finitely many coordinates and finitely many segments, a union bound yields
\[
\Pb_{\xi_0}^{\cK_{C,B}}\bigl(\|\widehat\mu_B-\sum_{m=1}^M (L_m/N)\mu_{f_m}^{R_m}\|_1>\eta/2
\ \bigm|\ G_B\bigr)\le B^{-3}
\]
for all sufficiently large \(B\).  By \eqref{eq:rounding},
\[
\left\|\sum_{m=1}^M (L_m/N)\mu_{f_m}^{R_m}-\mu_C\right\|_1\le \frac{2M}{N},
\]
which is at most \(\eta/4\) for large \(B\).  The reset windows occupy only \(r(B)\) out of
\(B\) stages, so their contribution to the empirical occupation is at most \(r(B)/B\le
\eta/4\) for large \(B\).  Combining these estimates on \(G_B\) gives
\[
\Pb_{\xi_0}^{\cK_{C,B}}\bigl(\|\widehat\mu_B-\mu_C\|_1>\eta\mid G_B\bigr)\le B^{-3}.
\]
Together with \(\Pb(G_B^c)\le B^{-3}\), this yields the stated
\(1-B^{-2}\) bound after enlarging \(B_0\) if necessary.  Publicity is clear from the
construction.
\end{proof}

\subsection{The local controller lemma for exit-capable nodes}
\label{subsec:local-controller}

We now lift \Cref{thm:mosaic} to augmented nodes.  Thanks to
\Cref{thm:graphVP3,prop:JC-extreme}, this is a purely affine matter.

\begin{theorem}[Local Controller Lemma]
\label{thm:local-controller}
Fix an exit-capable node \(C\), and let
\[
\bar\mu_C^{*}=J_C(\mu_C^{*})=(\mu_C^{*},L_C\mu_C^{*})\in \cM_C^{\mathrm{aug}}.
\]
For every \(\delta>0\) there exist a block length \(L\), constants \(K_C<\infty\) and
\(\rho_C\in(0,1)\), and a public finite-memory controller \(\cK_{C,\delta}\) such that,
from every entry phase-state \(\xi_0\in \Xi_C\),
\begin{enumerate}[label=\textnormal{(\roman*)},leftmargin=2.6em]
\item \(\cK_{C,\delta}\) depends only on public information and is deterministic once the
public memory is fixed;
\item if \(\widehat\mu_C^{\mathrm{aug}}\) denotes the expected augmented occupation vector
generated by the \(L\)-block of \(\cK_{C,\delta}\), then
\begin{equation}
\label{eq:local-controller-bound}
\bigl\|\widehat\mu_C^{\mathrm{aug}}-\bar\mu_C^{*}\bigr\|_1
\le
(1+\Lambda_C)\eta + K_C\rho_C^L
\le \delta,
\end{equation}
where \(\Lambda_C:=\|L_C\|_{1\to 1}\) and \(\eta=\delta/(2(1+\Lambda_C))\);
\item the required block length \(L\) depends on \(C\) and \(\delta\), but not on the
entry state.
\end{enumerate}
\end{theorem}

\begin{proof}
Fix \(\delta>0\), and set
\[
\eta:=\frac{\delta}{2(1+\Lambda_C)}.
\]
Apply \Cref{thm:mosaic} to \(\mu_C^{*}\) with accuracy \(\eta\).  Thus, for all large block
lengths \(L\), there is a public controller whose internal expected occupation
\(\widehat\mu_C\) satisfies
\[
\|\widehat\mu_C-\mu_C^{*}\|_1\le \eta
\]
up to an exponentially small finite-block remainder.  More precisely, the reset and
mixing estimates in the proof of \Cref{thm:mosaic} give constants \(K_C<\infty\) and
\(\rho_C\in(0,1)\) such that
\[
\|\widehat\mu_C-\mu_C^{*}\|_1\le \eta + K_C\rho_C^L.
\]

Now lift by the affine injection \(J_C\).  Since \(J_C(\nu)=(\nu,L_C\nu)\),
\[
\|J_C(\widehat\mu_C)-J_C(\mu_C^{*})\|_1
\le
\|\widehat\mu_C-\mu_C^{*}\|_1
+
\|L_C(\widehat\mu_C-\mu_C^{*})\|_1
\le
(1+\Lambda_C)\|\widehat\mu_C-\mu_C^{*}\|_1.
\]
Therefore
\[
\|J_C(\widehat\mu_C)-\bar\mu_C^{*}\|_1
\le
(1+\Lambda_C)\eta + (1+\Lambda_C)K_C\rho_C^L.
\]
Absorb the factor \(1+\Lambda_C\) into the constant \(K_C\).  Choosing \(L\) large enough
that \(K_C\rho_C^L\le \delta/2\) yields \eqref{eq:local-controller-bound}.  Since the
input controller from \Cref{thm:mosaic} is uniform over the entry state, the same is true
for the lifted controller.
\end{proof}

\begin{remark}
The statement is exactly the nodewise public realizability theorem that was missing in the
v3 and early v4 proof states.  The key point is that the augmented node is not a new
nonlinear object.  It is the graph of the linear map \(L_C\), so the same
Carath\'eodory mosaic realizes it.
\end{remark}

\subsection{Flux witness, exit routing, and the realization theorem}

Fix a relaxed fixed point
\[
x^*=(x_C^*)_{C\in \cT_\eta}\in \cA_\eta,
\qquad
x_C^*=\bigl(\mu_C^*,\beta_C^*,(v_C^*(b))_{b\in B_C},(g_{i,C}^*,h_{i,C}^*)_{i\in N}\bigr).
\]
By \Cref{thm:graphVP3}, every node package already carries its compatible exit masses.  No
additional witness needs to be chosen.

\begin{definition}
A sequence \((M_r)_{r\ge 1}\) of positive integers is a \emph{superexponential schedule}
if \(M_{r+1}\ge M_r^2\) for every \(r\).  Its \emph{frozen version at level \(R\)} is the
sequence \((M_r^{(R)})_{r\ge 1}\) defined by
\[
M_r^{(R)}=
\begin{cases}
M_r, & r<R,\\
M_R, & r\ge R.
\end{cases}
\]
We write \(M_*:=M_R\) for the tail block length.
\end{definition}

The frozen schedule is the device that makes the final strategy finite-memory while
retaining the summable-leakage property of the original growing blocks.

\begin{proposition}[Public exit routing]
\label{prop:routing}
Fix \(\eta>0\), \(\delta>0\), a relaxed fixed point \(x^*\in \cA_\eta\), and for each node
\(C\) a local controller \(\cK_{C,\delta}\) from \Cref{thm:local-controller}.  Then there
exists a public strategy profile \(\Sigma_{\eta,\delta}^{\mathrm{gb}}\) with the following
properties.
\begin{enumerate}[label=\textnormal{(\roman*)},leftmargin=2.6em]
\item The public memory stores the current node \(C\), the current block index \(r\), and
the current internal memory of the node controller.
\item On block \(r\) in node \(C\), the profile runs \(\cK_{C,\delta}\) for exactly
\(M_r\) stages.
\item If the block emits boundary state \(b\in B_C\), the next block begins in the child
node \(C[b]\); if no exit occurs, the next block remains in \(C\).
\item The routing update is publicly observable and deterministic conditional on the public
memory.
\end{enumerate}
\end{proposition}

\begin{proof}
The only issue is observability.  The local controller \(\cK_{C,\delta}\) is public, so its
internal state is part of the public memory.  The current node is public by induction on
blocks.  Boundary states are observed in the public state history, hence the identity of
the emitted \(b\in B_C\) is public.  The successor map \(b\mapsto C[b]\) is fixed in the
tree data, so the node update is deterministic.  Finally the block index \(r\) increments
deterministically.
\end{proof}

\begin{theorem}[Realization theorem]
\label{thm:realization}
Fix \(\eta>0\), a relaxed fixed point \(x^*\in \cA_\eta\), and \(\delta>0\).  Let
\((M_r)_{r\ge 1}\) be a superexponential schedule, and freeze it at level \(R\), with tail
length \(M_*=M_R\).  Then \Cref{prop:routing} yields a public finite-memory profile
\(\sigma_{\eta,\delta,M_*}\) and a constant \(K_\eta<\infty\) such that:
\begin{enumerate}[label=\textnormal{(R\arabic*)},leftmargin=2.8em]
\item for every complete block \(B\) played in node \(C\),
\begin{equation}
\label{eq:block-aug}
\bigl\|\widehat\mu_B^{\mathrm{aug}}-\bar\mu_C^*\bigr\|_1\le \delta;
\end{equation}
\item for every player \(i\), every complete block \(B\) played in node \(C\), and every
entry history to that block,
\begin{equation}
\label{eq:block-payoff-close}
\left|
\E\!\left[\frac1{|B|}\sum_{t\in B}u_i(s_t,a_t)\,\Bigm|\,\mathcal F_{\mathrm{in}(B)}\right]
-g_{i,C}^*
\right|
\le C_\eta \delta
\end{equation}
for some constant \(C_\eta<\infty\) depending only on the compact set \(\cA_\eta\) and the
game;
\item for every horizon \(T\),
\begin{equation}
\label{eq:realization-tail}
\left|
\gamma_i^T(s,\sigma_{\eta,\delta,M_*})
-
\frac1T\sum_{B\subseteq [1,T]} |B|\,g_{i,C(B)}^*
\right|
\le C_\eta\delta+\frac{K_\eta}{T},
\end{equation}
where the sum runs over complete blocks \(B\) contained in \([1,T]\), and \(C(B)\) denotes
the node of block \(B\).
\end{enumerate}
\end{theorem}

\begin{proof}
Property \eqref{eq:block-aug} is the content of \Cref{thm:local-controller}, applied to
the node visited by the block.  Because the frozen schedule uses only finitely many block
lengths, the resulting global profile has finite memory.

To prove \eqref{eq:block-payoff-close}, note that the block-average payoff is a continuous
affine functional of the augmented occupation vector \(\widehat\mu_B^{\mathrm{aug}}\): the
current payoff part is linear in the internal occupation coordinates, and the continuation
part is linear in the exit-flux coordinates.  Since \(x^*\) ranges over the compact set
\(\cA_\eta\), these affine functionals have a uniform Lipschitz constant \(C_\eta\).  The
gain identity in the definition of the local Nash set gives exactly \(g_{i,C}^*\) for the
target block vector \(\bar\mu_C^*\).  Combining this with \eqref{eq:block-aug} yields
\eqref{eq:block-payoff-close}.

For \eqref{eq:realization-tail}, decompose the horizon \(\{1,\dots,T\}\) into the complete
blocks contained in \([1,T]\) plus at most one initial incomplete fragment and one final
incomplete fragment.  The total contribution of the incomplete fragments is bounded by a
constant depending only on the frozen prelude and the tail block length \(M_*\); absorb
that constant into \(K_\eta\).  The complete blocks are controlled by summing
\eqref{eq:block-payoff-close} and dividing by \(T\).
\end{proof}

\section{Proof of Theorem A under condition \texorpdfstring{$(a)$}{(a)}}
\label{sec:proofA}

This section upgrades the realization theorem to a uniform equilibrium proof.

\subsection{Condition \texorpdfstring{$(a)$}{(a)} and quantitative irreducibility}

From this point onward we work in the repaired coordinates of Section~\ref{sec:flux}.  To
stress this, we sometimes write \(\widehat{\cA}_\eta\) for the same space \(\cA_\eta\).

For a node \(C\), a package \(x\in \cA_\eta\), a player \(i\), and a pure stationary
selector \(a_i:\Xi_C\to \bigsqcup_{\xi\in \Xi_C}A_i(s(\xi))\), let
\[
Q_{C,i}^{x,a_i}
\]
denote the unilateral node kernel obtained when the other players keep the local behavior
encoded by \(x_C\) and player \(i\) plays \(a_i\) inside node \(C\).  This kernel acts on
\(\Xi_C\).

\begin{condition}[Condition \textnormal{(a)}]
\label{cond:a}
For every \(\eta\in(0,1)\), every repaired package \(x\in \widehat{\cA}_\eta\), every node
\(C\in \cT_\eta\), every player \(i\in N\), and every pure stationary selector \(a_i\),
the kernel \(Q_{C,i}^{x,a_i}\) is irreducible on \(\Xi_C\).
\end{condition}

This condition is exactly the statement that no player can unilaterally break
communication inside any repaired node kernel.

\begin{lemma}
\label{lem:kappa}
Assume condition~\textnormal{(a)}.  For every \(\eta\in(0,1)\),
\[
\kappa_\eta
:=
\inf_{x\in \widehat{\cA}_\eta}\;
\min_{C\in \cT_\eta}\;
\min_{i\in N}\;
\min_{a_i}\;
\min_{\xi,\zeta\in \Xi_C}
\max_{1\le n\le |\Xi_C|-1}
\bigl(Q_{C,i}^{x,a_i}\bigr)^n(\xi,\zeta)
\]
satisfies \(\kappa_\eta>0\).
\end{lemma}

\begin{proof}
Fix a node \(C\), a player \(i\), and a selector \(a_i\).  By condition~(a), the kernel
\(Q_{C,i}^{x,a_i}\) is irreducible for every \(x\in \widehat{\cA}_\eta\).  For a finite
irreducible kernel on \(|\Xi_C|\) states, every pair \((\xi,\zeta)\) is connected by a path
of length at most \(|\Xi_C|-1\).  Hence for each such kernel,
\[
\min_{\xi,\zeta}\max_{1\le n\le |\Xi_C|-1} (Q_{C,i}^{x,a_i})^n(\xi,\zeta)>0.
\]
The map \(x\mapsto Q_{C,i}^{x,a_i}\) is continuous on the compact set
\(\widehat{\cA}_\eta\), and the displayed expression is a continuous function of the
kernel.  Therefore its minimum over \(\widehat{\cA}_\eta\) is attained and positive.
Finally there are only finitely many nodes, players, and selectors, so the minimum over
all of them is still positive.
\end{proof}

\subsection{Bias-span bound}

The importance of \(\kappa_\eta\) is that it yields a uniform bound on the span of all
local bias certificates.

\begin{lemma}[Expected hitting-time bound]
\label{lem:hitting}
Let \(Q\) be an irreducible Markov kernel on a finite set \(X\), and let
\[
\kappa(Q):=\min_{x,y\in X}\max_{1\le n\le |X|-1}Q^n(x,y).
\]
Then for all \(x,y\in X\),
\[
\E_x[\tau_y]\le \frac{|X|-1}{\kappa(Q)},
\qquad
\tau_y:=\inf\{t\ge 0:X_t=y\}.
\]
\end{lemma}

\begin{proof}
Fix \(x,y\in X\), and set \(m:=|X|-1\).  By definition of \(\kappa(Q)\),
\[
\Pb_x(\tau_y\le m)\ge \kappa(Q).
\]
Applying the Markov property after each block of \(m\) steps gives
\[
\Pb_x(\tau_y>km)\le (1-\kappa(Q))^k
\qquad (k\ge 0).
\]
Therefore
\[
\E_x[\tau_y]
=
\sum_{t\ge 0}\Pb_x(\tau_y>t)
\le
m\sum_{k\ge 0}\Pb_x(\tau_y>km)
\le
m\sum_{k\ge 0}(1-\kappa(Q))^k
=
\frac{m}{\kappa(Q)}.
\]
\end{proof}

\begin{proposition}[Bias-span bound]
\label{prop:bias-span}
Assume condition~\textnormal{(a)}.  Fix \(\eta\in(0,1)\), a node \(C\in \cT_\eta\), and a
player \(i\in N\).  Let \((g_{i,C}^*,h_{i,C}^*)\) be a local Bellman certificate at the
fixed point \(x^*\in \cA_\eta\), normalized by
\(\min_{\xi\in \Xi_C} h_{i,C}^*(\xi)=0\).  Then
\[
\spn(h_{i,C}^*)
\le
H_\eta
:=
\frac{\max_D|\Xi_D|-1}{\kappa_\eta}.
\]
\end{proposition}

\begin{proof}
Fix \(\xi\in \Xi_C\), and choose \(\zeta\in \Xi_C\) with
\(h_{i,C}^*(\zeta)=\min_{\xi'} h_{i,C}^*(\xi')=0\).  Let \(a_i^*\) be any pure stationary
selector that attains the maximizing action in the Bellman inequality at each state.
Because the local Bellman certificate is exact, along the Markov chain with kernel
\(Q:=Q_{C,i}^{x^*,a_i^*}\) one has
\[
u_i(X_t,A_t)+\E[h_{i,C}^*(X_{t+1})\mid \mathcal F_t]-h_{i,C}^*(X_t)\le g_{i,C}^*
\qquad\text{a.s.}
\]
Since \(u_i\in[0,1]\) and \(g_{i,C}^*\in[0,1]\), it follows that
\[
\E[h_{i,C}^*(X_{t+1})-h_{i,C}^*(X_t)\mid \mathcal F_t]\le 1.
\]
Stopping at \(\tau_\zeta\) and summing gives
\[
h_{i,C}^*(\xi)-h_{i,C}^*(\zeta)\le \E_\xi[\tau_\zeta].
\]
By \Cref{lem:hitting},
\[
\E_\xi[\tau_\zeta]\le \frac{|\Xi_C|-1}{\kappa(Q)}\le \frac{\max_D|\Xi_D|-1}{\kappa_\eta},
\]
because \(|\Xi_C|\le |S|\) and \(\kappa(Q)\ge \kappa_\eta\) by definition of \(\kappa_\eta\).
Since \(h_{i,C}^*(\zeta)=0\), this proves
\[
\max_{\xi} h_{i,C}^*(\xi)\le \frac{\max_D|\Xi_D|-1}{\kappa_\eta}.
\]
The minimum is \(0\), so the span is at most the same quantity.
\end{proof}

\subsection{Deviation-value continuity under condition~(a)}

The following lemma provides the key quantitative bridge from compliant occupation closeness to deviation-value control. It is the central ingredient that makes the blockwise Bellman certificate rigorous.

\begin{lemma}[Deviation-value Lipschitz bound]\label{lem:dev-value-lipschitz}
Fix a node \(C\) and a player \(i\). Assume condition~(a), and fix \(\eta>0\). Let \(\mu,\widehat\mu\in\mathcal P_C\) be two admissible occupation vectors for node \(C\). Denote by \(V^{\mathrm{dev}}_{i,C}(\mu)\) and \(V^{\mathrm{dev}}_{i,C}(\widehat\mu)\) the optimal average deviation values for player \(i\) in the unilateral deviation MDP induced at node \(C\) by the environments \(\mu\) and \(\widehat\mu\), respectively. Then
\[
\bigl|V^{\mathrm{dev}}_{i,C}(\widehat\mu)-V^{\mathrm{dev}}_{i,C}(\mu)\bigr|
\le \operatorname{sp}(h_{i,C,\mu})\,\|\widehat\mu-\mu\|_1,
\]
where \(h_{i,C,\mu}\) is any bias function for the deviation MDP associated with \(\mu\). Moreover, under condition~(a),
\[
\operatorname{sp}(h_{i,C,\mu})\le \frac{|\Xi_C|-1}{\kappa_\eta},
\]
and therefore
\[
\bigl|V^{\mathrm{dev}}_{i,C}(\widehat\mu)-V^{\mathrm{dev}}_{i,C}(\mu)\bigr|
\le \frac{\max_{D}|\Xi_D|-1}{\kappa_\eta}\,\|\widehat\mu-\mu\|_1
=H_\eta\,\|\widehat\mu-\mu\|_1.
\]
\end{lemma}

\begin{proof}
Fix \(C\) and \(i\). For each admissible occupation vector \(\nu\in\mathcal P_C\), the unilateral optimization problem of player \(i\) against the compliant environment determined by \(\nu\) is a finite average-reward Markov decision process on \(\Xi_C\). Let \(Q_\nu\) denote its transition kernel, \(r_\nu\) its one-stage reward vector, and \(g_\nu=V^{\mathrm{dev}}_{i,C}(\nu)\) its optimal average reward. Let \(h_\nu\) be a bias function satisfying the average-reward optimality equation
\begin{equation}\label{eq:bellman-dev-lipschitz}
g_\nu \mathbf 1 + h_\nu = T_\nu h_\nu,
\end{equation}
where \(T_\nu\) is the dynamic programming operator
\[
(T_\nu f)(\xi)=\max_{a_i\in A_i(\xi)}\Bigl\{r_\nu(\xi,a_i)+\sum_{\xi'}Q_\nu(\xi' \mid \xi,a_i)f(\xi')\Bigr\}.
\]

Apply \eqref{eq:bellman-dev-lipschitz} with \(\nu=\mu\). Evaluating the \(\widehat\mu\)-operator at \(h_\mu\) and using \(|\max_u \alpha_u-\max_u \beta_u|\le \max_u |\alpha_u-\beta_u|\), we get
\[
\|T_{\widehat\mu}h_\mu-T_\mu h_\mu\|_\infty
\le
\operatorname{sp}(h_\mu)\,\|\widehat\mu-\mu\|_1.
\]
The bound arises because the one-stage reward and transition data depend affinely on the occupation vector: the reward perturbation is at most \(\|\widehat\mu-\mu\|_1\), and the transition perturbation satisfies \(|(Q_{\widehat\mu}-Q_\mu)(\xi,a_i)h_\mu| \le \operatorname{sp}(h_\mu)\,\|Q_{\widehat\mu}(\xi,a_i)-Q_\mu(\xi,a_i)\|_{\mathrm{TV}} \le \operatorname{sp}(h_\mu)\,\|\widehat\mu-\mu\|_1\).

Now let \(\pi\) be an optimal stationary deviation policy for the MDP associated with \(\widehat\mu\). Testing with \(h_\mu\) gives \(g_{\widehat\mu} \le g_\mu+\operatorname{sp}(h_\mu)\,\|\widehat\mu-\mu\|_1\). Interchanging roles yields the reverse inequality, so
\[
|g_{\widehat\mu}-g_\mu|
\le
\max\{\operatorname{sp}(h_\mu),\operatorname{sp}(h_{\widehat\mu})\}\,\|\widehat\mu-\mu\|_1.
\]

It remains to bound the bias span uniformly. Under condition~(a), every admissible unilateral deviation kernel on node \(C\) retains connector probabilities bounded below by \(\kappa_\eta>0\). Therefore the deviation MDP is communicating on \(\Xi_C\), with hitting times bounded by \(|\Xi_C|-1)/\kappa_\eta\). For communicating finite average-reward MDPs, the span of a bias function is bounded by the maximal hitting time, giving \(\operatorname{sp}(h_\nu)\le (|\Xi_C|-1)/\kappa_\eta\) for every admissible \(\nu\).
\end{proof}

\subsection{Blockwise Bellman certificate}

Fix \(\eta\in(0,1)\), choose a relaxed fixed point \(x^*\in \cA_\eta\), and let
\(\sigma_{\eta,\delta,M_*}\) be the realized profile from \Cref{thm:realization}.  Write
\[
H_\eta:=\frac{\max_D|\Xi_D|-1}{\kappa_\eta}.
\]
By \Cref{prop:bias-span}, each local bias certificate \(h_{i,C}^*\) can be chosen with
span at most \(H_\eta\).

\begin{lemma}
\label{lem:block-bellman}
There exists a constant \(C_\eta<\infty\) such that for every \(\delta>0\), every complete
block \(B\) of the realized profile played in node \(C\), every player \(i\), and every
unilateral deviation \(\tau_i\),
\begin{equation}
\label{eq:block-bellman}
\E\!\left[
\frac1{|B|}\sum_{t\in B} u_i^t(\tau_i,\sigma_{-i})
\,\Bigm|\, \mathcal F_{\mathrm{in}(B)}
\right]
\le
g_{i,C}^*+\eta+C_\eta\delta+\frac{2H_\eta}{|B|}.
\end{equation}
\end{lemma}

\begin{proof}
Fix a block \(B=\{t_0,\dots,t_1-1\}\) of length \(M:=|B|\), played in node \(C\).  The
exact local Bellman inequality at the fixed point says that for every phase-state
\(\xi\) and every action \(a_i\in A_i(s(\xi))\),
\[
r_{i,C}^*(\xi,a_i)+\sum_{\xi'}Q_{C,i}^*(\xi' \mid \xi,a_i)\,h_{i,C}^*(\xi')
\le
g_{i,C}^*+h_{i,C}^*(\xi)+\eta.
\]
Here \(r_{i,C}^*\) and \(Q_{C,i}^*\) are the reward and transition coefficients induced by
the fixed-point package \(x_C^*\).  Because the realized block occupation \(\widehat\mu\) is within
\(\delta\) of \(\bar\mu_C^*\) in the sense of \eqref{eq:block-aug}, we transfer the certificate \(h_{i,C}^*\) to the realized block as follows. The reward and transition data depend \emph{affinely} on the occupation vector: for each state-action pair \((\xi,a_i)\),
\[
\|r_{\widehat\mu}(\xi,a_i)-r_{\mu_C^*}(\xi,a_i)\|\le \|\widehat\mu-\mu_C^*\|_1\le\delta,
\]
and
\[
\|Q_{\widehat\mu}(\cdot\mid\xi,a_i)-Q_{\mu_C^*}(\cdot\mid\xi,a_i)\|_{\mathrm{TV}}\le \|\widehat\mu-\mu_C^*\|_1\le\delta.
\]
These are the \emph{coefficient-level} perturbation bounds (not merely value-level), following from the affine dependence of the block dynamics on the occupation measure. Combined with \(\operatorname{sp}(h_{i,C}^*)\le H_\eta\) from \Cref{lem:dev-value-lipschitz}, we get for every \((\xi,a_i)\):
\[
\bigl|r_{\widehat\mu}(\xi,a_i)+Q_{\widehat\mu}(\xi,a_i)h_{i,C}^* - r_{\mu_C^*}(\xi,a_i)-Q_{\mu_C^*}(\xi,a_i)h_{i,C}^*\bigr|
\le \delta+H_\eta\,\delta = (1+H_\eta)\delta.
\]
Set \(C_\eta:=1+H_\eta\). The exact Bellman inequality at the fixed point therefore transfers to the realized block: for every action \(a_i\),
\[
r_{\widehat\mu}(\xi,a_i)+Q_{\widehat\mu}(\xi,a_i)h_{i,C}^*
\le g_{i,C}^*+h_{i,C}^*(\xi)+\eta+C_\eta\delta.
\]
Therefore, along the deviating play inside the block,
\[
\E\!\left[
u_i^t(\tau_i,\sigma_{-i})
+h_{i,C}^*(X_{t+1})-h_{i,C}^*(X_t)
\,\Bigm|\, \mathcal F_t
\right]
\le
g_{i,C}^*+\eta+C_\eta\delta
\]
for every \(t\in B\).  Summing from \(t_0\) to \(t_1-1\) and telescoping yields
\[
\E\!\left[
\sum_{t\in B}u_i^t(\tau_i,\sigma_{-i})
\,\Bigm|\, \mathcal F_{\mathrm{in}(B)}
\right]
\le
M(g_{i,C}^*+\eta+C_\eta\delta)
+
\E\!\left[h_{i,C}^*(X_{t_0})-h_{i,C}^*(X_{t_1})\mid \mathcal F_{\mathrm{in}(B)}\right].
\]
The terminal bias difference is bounded by \(\spn(h_{i,C}^*)\le H_\eta\) at each end, so
its absolute value is at most \(2H_\eta\).  Dividing by \(M\) gives
\eqref{eq:block-bellman}.
\end{proof}

\begin{corollary}
\label{cor:block-regret}
For every complete block \(B\) played in node \(C\),
\[
\E\!\left[
\frac1{|B|}\sum_{t\in B} u_i^t(\tau_i,\sigma_{-i})
-
\frac1{|B|}\sum_{t\in B} u_i^t(\sigma)
\,\Bigm|\, \mathcal F_{\mathrm{in}(B)}
\right]
\le
\eta + C_\eta\delta + \frac{2H_\eta}{|B|}.
\]
After possibly increasing \(C_\eta\), the same constant as in \Cref{lem:block-bellman}
may be used.
\end{corollary}

\begin{proof}
By \Cref{lem:block-bellman}, the deviating block average is at most
\(g_{i,C}^*+\eta+C_\eta\delta+2H_\eta/|B|\).  By \eqref{eq:block-payoff-close}, the block
average under the realized profile differs from \(g_{i,C}^*\) by at most \(C_\eta\delta\).
Renaming the constant finishes the proof.
\end{proof}

\subsection{Proof of Theorem A}

We now assemble the estimates.

\begin{proof}[Proof of \Cref{thm:A}]
Fix \(\eps>0\).  Choose parameters in the order
\[
\eta,\qquad \delta,\qquad M_*,\qquad T_0.
\]

\smallskip
\noindent\textbf{Step 1: the relaxed fixed point.}
Choose \(\eta>0\) so small that \(\eta\le \eps/4\).  By \Cref{thm:Kakutani}, there exists
a relaxed fixed point \(x^*\in \cA_\eta\).

\smallskip
\noindent\textbf{Step 2: irreducibility and span.}
Because condition~(a) holds, \Cref{lem:kappa} gives \(\kappa_\eta>0\).  Define
\[
H_\eta:=\frac{\max_D|\Xi_D|-1}{\kappa_\eta}.
\]
By \Cref{prop:bias-span}, each local bias certificate at \(x^*\) can be chosen with span
at most \(H_\eta\).

\smallskip
\noindent\textbf{Step 3: realization.}
Choose \(\delta>0\) such that \(C_\eta\delta\le \eps/4\).  Fix a superexponential schedule,
freeze it at a sufficiently late level \(R\), and let \(M_*=M_R\).  Apply
\Cref{thm:realization} to obtain the public finite-memory profile
\[
\sigma^\eps:=\sigma_{\eta,\delta,M_*}.
\]

\smallskip
\noindent\textbf{Step 4: switching loss.}
Increase \(R\) if necessary so that
\[
\frac{2H_\eta}{M_*}\le \frac{\eps}{4}.
\]
Because the frozen schedule uses only finitely many pre-tail blocks, the realization
theorem yields a constant \(K_\eta<\infty\) such that the incomplete blocks and frozen
prelude contribute at most \(K_\eta/T\).

\smallskip
\noindent\textbf{Step 5: the horizon threshold.}
Choose \(T_0\) so large that
\[
\frac{K_\eta}{T_0}\le \frac{\eps}{4}.
\]

Now fix \(T\ge T_0\), an initial state \(s\), a player \(i\), and a unilateral deviation
\(\tau_i\).  Partition \(\{1,\dots,T\}\) into the complete blocks contained in \([1,T]\)
plus the incomplete initial and terminal fragments.  Summing
\Cref{cor:block-regret} over the complete blocks and using that each complete block has
length at least \(M_*\) yields
\[
\gamma_i^T\bigl(s,(\tau_i,\sigma^\eps_{-i})\bigr)
-
\gamma_i^T(s,\sigma^\eps)
\le
\eta + C_\eta\delta + \frac{2H_\eta}{M_*} + \frac{K_\eta}{T}.
\]
By the choices above, the right-hand side is at most
\[
\frac{\eps}{4}+\frac{\eps}{4}+\frac{\eps}{4}+\frac{\eps}{4}=\eps.
\]
Therefore
\[
\gamma_i^T(s,\sigma^\eps)
\ge
\gamma_i^T\bigl(s,(\tau_i,\sigma^\eps_{-i})\bigr)-\eps.
\]
Since \(s\), \(i\), \(\tau_i\), and \(T\ge T_0\) were arbitrary, \(\sigma^\eps\) is a
uniform \(\eps\)-equilibrium.  The construction is public finite-memory by
\Cref{thm:realization}.  This proves \Cref{thm:A}.
\end{proof}

\begin{remark}
The proof uses condition~(a) in one place only: the lower bound \(\kappa_\eta>0\) needed
for the span estimate.  Everything else---compactness, graph-VP3, Carath\'eodory mosaics,
exit-tag augmentation, and block realization---is unconditional.
\end{remark}

\section{Discussion}
\label{sec:discussion}

\subsection{What condition \texorpdfstring{$(a)$}{(a)} buys}

The result of this paper is deliberately sharp about where the argument needs its extra
assumption.  The compact asymptotic state space, the flux repair, the graph theorem, and
the nodewise realization machinery are unconditional.  The only missing bridge in the
general case is the possibility that a player may disconnect a node by a unilateral local
selector.  Condition~(a) rules this out, and once it is ruled out the Bellman certificate
has a uniformly bounded span.  That bounded span is the exact quantity needed to control
the block-boundary telescoping terms.

The theorem should therefore be read as saying that the compact-tree route does not fail
because of topological or realization issues.  Those issues are solved here.  The route
fails, if it fails at all, at the communication level.

\subsection{Why graph-VP3 matters}

The v3 draft carried a separate compatibility condition between internal occupation and
exit witnesses.  The present paper shows that this was the wrong formulation.  The product
condition is too large: it asks for an entire rectangle of combinations that are not
jointly feasible.  The correct feasible set is the graph of the linear exit map.  This is
why the local controller lemma is possible: once the node package is regarded as
\((\mu,L_C\mu)\), every Carath\'eodory decomposition of \(\mu\) automatically lifts.

This observation is conceptually important beyond the present theorem.  It suggests that in
other asymptotic constructions one should expect boundary data to live on a graph, not in
a product.  Normalized exit shares look natural, but they keep artificial directions alive
when the total exit mass vanishes.

\subsection{Relation to the full Neyman conjecture}

Theorem~A does not settle Neyman's conjecture.  The full conjecture asks for uniform
equilibria without any nodewise irreducibility hypothesis.  Still, the current proof
isolates the precise obstruction encountered by the compact asymptotic route:
communication can degenerate under unilateral local deviations even when the relaxed fixed
point itself is well behaved.

There are at least three natural next questions.

\begin{enumerate}[label=\textnormal{(\arabic*)},leftmargin=2.4em]
\item Can condition~(a) be deduced from softer structural hypotheses, perhaps after a
further refinement of the node decomposition?
\item Is there a replacement for the bias-span argument that works without irreducibility,
for example by carrying an additional singular part in the local potential?
\item Can one construct a counterexample in which the compact tree fixed point exists and
is realizable, but the lack of unilateral communication produces unavoidable regret spikes?
\end{enumerate}

The present paper does not answer these questions.  It does, however, turn them into
honest mathematical questions rather than artifacts of an incomplete implementation route.

\appendix

\section{The local average-reward linear programs}
\label{app:local-LP}

This appendix records the finite-dimensional reason the local relaxed Nash sets used in
Section~\ref{sec:tree} are nonempty, compact, and closed-graph.

Fix \(\eta\in(0,1)\), a node \(C\), a player \(i\), and an environment \(x\in \cA_\eta\).
Because all underlying sets are finite, the unilateral control problem of player \(i\)
inside node \(C\) is a finite average-reward Markov decision problem once the other
players' local behavior and all continuation values are frozen according to \(x\).  Its
primal variables are occupation measures of the unilateral control problem, and its dual
variables are the gain and bias pair \((g,h)\).

\begin{proposition}
\label{prop:local-LP}
For fixed \(C,i,\eta\), there exist matrices and vectors depending continuously on the
environment \(x\in\cA_\eta\),
\[
A_{C,i}(x),\quad b_{C,i}(x),\quad F_C(x),
\]
such that the exact local certificate set \(\widetilde{\cN}_C(x)\) is the projection of
the compact feasible set of the finite linear system
\[
A_{C,i}(x)z\le b_{C,i}(x),
\qquad
F_C(x)z=0,
\]
together with the compact box constraints already built into \(\cA_\eta\).
Consequently \(\widetilde{\cN}_C(x)\) is nonempty and compact, and the map
\(x\mapsto \widetilde{\cN}_C(x)\) has closed graph.
\end{proposition}

\begin{proof}
The unilateral control problem is a finite average-reward MDP.  The standard occupation
measure formulation gives a primal linear program in the unilateral occupation variables.
The standard Bellman inequalities give the dual linear program in \((g,h)\).  Strong
duality holds because both programs are feasible and finite-dimensional.  The exact local
certificate set consists of primal-feasible and dual-feasible points with equal objective
value, plus the continuation-matching constraints.  All coefficients depend continuously on
the environment \(x\): rewards depend affinely on the continuation values, and transition
coefficients depend continuously on the local data stored in \(x_C\).  Since the node
package space already imposes compact box constraints, the feasible set is compact.  The
projection of a compact set is compact, proving nonemptiness and compactness of
\(\widetilde{\cN}_C(x)\).

For the graph property, let \(x^m\to x\) and \(y^m\to y\) with
\(y^m\in \widetilde{\cN}_C(x^m)\).  Lift each \(y^m\) to a feasible primal-dual point
\(z^m\) of the linear system.  By compactness, after passing to a subsequence we may
assume \(z^m\to z\).  Continuity of the coefficients gives
\(A_{C,i}(x)z\le b_{C,i}(x)\) and \(F_C(x)z=0\), so \(z\) is feasible for the limiting
system.  Projecting back yields \(y\in \widetilde{\cN}_C(x)\).
\end{proof}

\begin{remark}
The proposition is intentionally stated abstractly.  Writing the full matrices would
obscure the main line of the paper and add no conceptual content: the state and action
sets are finite, so the local Bellman system is an ordinary finite linear program.
\end{remark}

\section{A coarse universal bias box}
\label{app:bias-box}

The compactness proof in Section~\ref{sec:tree} used a universal bound \(B_\eta^0\) on the
local bias variables.  Any sufficiently large finite bound would do.  One convenient way
to produce such a bound is to normalize the dual LP from Appendix~\ref{app:local-LP} by
fixing \(\min_\xi h(\xi)=0\) and then selecting, for each local environment, an extreme
optimal solution.  The set of all environments is compact, and the set of extreme optimal
dual solutions of the finite LP varies upper hemicontinuously, so the sup norm of the
selected bias vectors is bounded on \(\cA_\eta\).  We denote any such bound by
\(B_\eta^0\).

The much sharper bound used in the proof of Theorem~A is the quantitative estimate
\(H_\eta=(\max_D|\Xi_D|-1)/\kappa_\eta\).  That sharper estimate requires condition~(a); the present
appendix does not.


\section{A concrete toy failure of product-VP3}
\label{app:toy-vp3}

This appendix gives a concrete finite-dimensional toy model illustrating the distinction
between the false product formulation and the correct graph formulation of VP3.  The point
is not that the toy model is itself a full stochastic game node; rather, it isolates the
only algebraic issue that mattered in the proof route.

Consider the line segment
\[
\cP:=
\{\mu^t:=t\mu^1+(1-t)\mu^0:0\le t\le 1\}
\subset \R^2,
\]
where
\[
\mu^0=(1,0),\qquad \mu^1=(0,1).
\]
Define the linear exit map
\[
L:\R^2\to \R^2,\qquad
L(x_1,x_2)=(x_1,x_2).
\]
The feasible joint set is therefore
\[
\mathcal J
=
\{(\mu,L\mu):\mu\in \cP\}
=
\{(\mu^t,\mu^t):0\le t\le 1\}.
\]
Its projection onto the occupation coordinates is \(\cP\), and its projection onto the exit
coordinates is again \(L(\cP)=\cP\).

\begin{proposition}
\label{prop:toy-rectangle}
In the toy model above, the product set
\[
\cP\times L(\cP)
\]
strictly contains the feasible joint set \(\mathcal J\).
\end{proposition}

\begin{proof}
The points
\[
(\mu^0,\mu^0),\qquad (\mu^1,\mu^1)
\]
belong to \(\mathcal J\).  So does every convex combination
\[
(\mu^t,\mu^t)=t(\mu^1,\mu^1)+(1-t)(\mu^0,\mu^0),\qquad 0\le t\le 1.
\]
But the mixed corner
\[
(\mu^0,\mu^1)
\]
belongs to \(\cP\times L(\cP)\) and does not belong to \(\mathcal J\), because every point
of \(\mathcal J\) has identical occupation and exit coordinates.  Hence
\(\mathcal J\subsetneq \cP\times L(\cP)\).
\end{proof}

The same geometry persists in the genuine node problem.  There the occupation polytope is
typically much larger than a segment, and the linear exit map is more complicated than the
identity, but the feasible joint set is still a graph.  Product-VP3 asks for the whole
rectangle and is therefore generically false.

A slightly richer picture, closer to the two-state gadget from the proof passes, is the
following.  Suppose
\[
L(\mu^t)=(t,1-t),
\]
so the exit flux can be viewed as a split between two boundary states.  Then the feasible
joint set is
\[
\{(\mu^t,(t,1-t)):0\le t\le 1\},
\]
again a graph over \(\cP\).  The product set would additionally contain mismatched points
such as \((\mu^0,(1/2,1/2))\) and \((\mu^{1/2},(1,0))\), none of which can arise from a
single controller.  This is the algebraic content of graph-VP3.

\section{Aggregating block estimates}
\label{app:block-aggregation}

This appendix isolates the deterministic bookkeeping behind the regret bound in
Section~\ref{sec:proofA}.

Let a horizon \(T\) be partitioned into pairwise disjoint intervals
\[
I^{\mathrm{pre}},\ B_1,\dots,B_m,\ I^{\mathrm{post}},
\]
where \(B_1,\dots,B_m\) are the complete blocks contained in \(\{1,\dots,T\}\), while
\(I^{\mathrm{pre}}\) and \(I^{\mathrm{post}}\) are the initial and terminal incomplete
fragments.  Either incomplete fragment may be empty.

\begin{lemma}
\label{lem:aggregation}
Assume that there exist constants \(a\ge 0\), \(H\ge 0\), \(M_*\ge 1\), and \(K<\infty\)
such that:
\begin{enumerate}[label=\textnormal{(\roman*)},leftmargin=2.6em]
\item every complete block satisfies \(|B_r|\ge M_*\);
\item for every \(r\),
\[
\E\!\left[
\sum_{t\in B_r} \Delta_t \,\Bigm|\, \mathcal F_{\mathrm{in}(B_r)}
\right]
\le
a|B_r|+2H,
\]
where \(\Delta_t\) is the stagewise regret increment;
\item the total contribution of the incomplete fragments satisfies
\[
\E\!\left[\sum_{t\in I^{\mathrm{pre}}\cup I^{\mathrm{post}}}\Delta_t\right]\le K.
\]
\end{enumerate}
Then
\[
\frac1T\,\E\!\left[\sum_{t=1}^T \Delta_t\right]
\le
a+\frac{2H}{M_*}+\frac{K}{T}.
\]
\end{lemma}

\begin{proof}
Take expectations in the blockwise estimate and sum over \(r=1,\dots,m\):
\[
\E\!\left[\sum_{r=1}^m\sum_{t\in B_r}\Delta_t\right]
\le
a\sum_{r=1}^m |B_r| + 2Hm.
\]
Since every block has length at least \(M_*\),
\[
m\le \frac{1}{M_*}\sum_{r=1}^m |B_r|.
\]
Hence
\[
\E\!\left[\sum_{r=1}^m\sum_{t\in B_r}\Delta_t\right]
\le
\left(a+\frac{2H}{M_*}\right)\sum_{r=1}^m |B_r|.
\]
Add the incomplete fragments and use assumption~(iii):
\[
\E\!\left[\sum_{t=1}^T\Delta_t\right]
\le
\left(a+\frac{2H}{M_*}\right)\sum_{r=1}^m |B_r| + K
\le
\left(a+\frac{2H}{M_*}\right)T + K.
\]
Divide by \(T\).
\end{proof}

In the proof of Theorem~A one applies \Cref{lem:aggregation} with
\[
a=\eta+C_\eta\delta,
\]
while the \(2H_\eta\) term comes from the blockwise telescope of the local bias function.

\begin{lemma}
\label{lem:frozen-memory}
A frozen superexponential schedule uses only finite memory.
\end{lemma}

\begin{proof}
Fix a freezing level \(R\).  The schedule uses only the finitely many lengths
\[
M_1,\dots,M_R.
\]
To implement it one only needs to store: the current node of the tree, the current stage
inside the current block, the local public memory of the node controller, and the current
block index truncated at \(R\).  Once the index reaches \(R\), it stays there forever.
All these coordinates range over finite sets, so the resulting profile is finite-memory.
\end{proof}

\begin{remark}
The original growing-block route uses the true superexponential schedule without freezing.
That profile is not finite-memory.  Freezing is the standard way to keep the summable
leakage estimates while producing an actual finite automaton.
\end{remark}


\begin{thebibliography}{99}

\bibitem{AshkenaziGolanKrasikovRainerSolan2024}
G.~Ashkenazi-Golan, I.~Krasikov, C.~Rainer, and E.~Solan.
\newblock Absorption paths and equilibria in quitting games.
\newblock \emph{Mathematical Programming}, 203(1--2):735--762, 2024.

\bibitem{BewleyKohlberg1976}
T.~Bewley and E.~Kohlberg.
\newblock The asymptotic theory of stochastic games.
\newblock \emph{Mathematics of Operations Research}, 1(3):197--208, 1976.

\bibitem{MertensNeyman1981}
J.-F.~Mertens and A.~Neyman.
\newblock Stochastic games.
\newblock \emph{International Journal of Game Theory}, 10(2):53--66, 1981.

\bibitem{MertensSorinZamir2003}
J.-F.~Mertens, S.~Sorin, and S.~Zamir.
\newblock Stochastic games.
\newblock In R.~Aumann and S.~Hart, editors, \emph{Handbook of Game Theory with
Economic Applications}, volume~3, chapter~47, pages 1811--1832. Elsevier, 2003.

\bibitem{Shapley1953}
L.~S.~Shapley.
\newblock Stochastic games.
\newblock \emph{Proceedings of the National Academy of Sciences of the United States of
America}, 39(10):1095--1100, 1953.

\bibitem{Solan1999}
E.~Solan.
\newblock Three-player absorbing games.
\newblock \emph{Mathematics of Operations Research}, 24(3):669--698, 1999.

\bibitem{Solan2003}
E.~Solan.
\newblock Uniform equilibrium: More than two players.
\newblock In A.~Neyman and S.~Sorin, editors, \emph{Stochastic Games and Applications},
NATO Science Series~570, pages 285--294. Springer, Dordrecht, 2003.

\bibitem{SolanVieille2025}
E.~Solan and N.~Vieille.
\newblock Undiscounted equilibrium in positive recursive absorbing games with
non-rectangular absorption structure.
\newblock \emph{arXiv preprint arXiv:2512.04306}, 2025.

\bibitem{Vieille2000I}
N.~Vieille.
\newblock Equilibrium in two-person stochastic games {I}: A reduction.
\newblock \emph{Israel Journal of Mathematics}, 119:55--91, 2000.

\bibitem{Vieille2000II}
N.~Vieille.
\newblock Equilibrium in two-person stochastic games {II}: The case of recursive games.
\newblock \emph{Israel Journal of Mathematics}, 119:93--126, 2000.

\end{thebibliography}

\end{document}
